\section{Conclusion}
In this project, we extended the railway scheduling problem within the Flatland framework by modeling energy consumption using Answer Set Programming (ASP). Our results demonstrate that the ASP encoding accurately captures physical realities at the local level: energy costs dynamically scale based on vehicle mass, power type, and specific kinematic states.

However, our current ASP encoding exhibits limitations regarding computational scalability. The introduction of complex energy penalties expands the search space significantly, particularly because the solver must evaluate numerous kinematic combinations across extended time horizons.

We consider several extensions and modifications for future work based on our findings:
\begin{itemize}
    \item \textbf{Algorithmic Improvements:} To manage the exponential growth of the search space on larger grids, future research could employ different solving techniques. Progressively extending the planning horizon could help approximate optimal energy solutions more efficiently without grounding the entire time horizon at once.
    \item \textbf{Environmental Physics:} With dynamic acceleration and momentum mechanics now established, future iterations could take track elevation and gradients/slopes into account, adding another layer of realism to how trains consume power on inclines or recover energy on declines.
    \item \textbf{Alternative Optimization Strategies:} Adjusting or combining the minimize directives and/or the heuristics could allow the solver to actively evaluate and choose longer, more energy-efficient detours rather than strictly defaulting to the fastest geometric route. Another strategy may also be to use Incremental Solving: This strategy was employed by Osswald and Leiber previously to minimize the sum of arrival times of all their agents for multiple environments, which was successful for a majority of their testing environments. However, as they mention, the main limitation of using Incremental Solving is that it only reduces the time horizon and, therefore, may not be suitable for all optimization metrics (\cite{osswald_railway_2026}).
\end{itemize}

Overall, we provide a robust proof of concept that ASP can effectively model physical constraints in railway logistics, laying the groundwork for future research into multi-objective routing.